\section{Fundamentação Teórica}
\label{sec:background}

O preenchimento de lacunas em imagens de satélite pode ser formulado como um problema de reconstrução condicionada por observações parciais.
Parte dos pixels permanece disponível e parte da informação é removida pela máscara de lacuna.
O objetivo do algoritmo é inferir os valores ausentes a partir do contexto espacial e espectral ainda observável.
Em imagens multiespectrais, esse problema é particularmente sensível porque a reconstrução não precisa apenas parecer plausível visualmente.
Ela precisa preservar relações radiométricas e espectrais que serão reutilizadas em tarefas analíticas posteriores.

Sob essa perspectiva, qualidade de reconstrução não é um conceito unidimensional.
Uma reconstrução pode apresentar baixo erro quadrático médio e ainda assim distorcer bordas, suavizar estruturas importantes ou alterar a relação entre bandas.
Da mesma forma, um método pode produzir boa média global e falhar sistematicamente em cenários mais complexos ou mais ruidosos.
Por isso, a análise do problema exige mais de uma métrica e também exige variáveis descritivas que ajudem a caracterizar a dificuldade da cena reconstruída.

A entropia local cumpre justamente esse papel.
Quando calculada sobre a vizinhança da lacuna, ela atua como um resumo da diversidade tonal e da organização espacial do entorno disponível para a inferência.
Regiões homogêneas tendem a exibir baixa entropia.
Regiões com textura, mistura de materiais, bordas e transições rápidas tendem a exibir entropia mais elevada.
Essa medida não substitui descrição semântica da cena e tampouco explica sozinha todo o comportamento do erro, mas oferece um eixo quantitativo útil para diferenciar regimes de dificuldade.

O uso de múltiplas janelas de entropia acrescenta uma camada importante à interpretação.
Uma janela pequena enfatiza detalhe fino e variações muito locais.
Uma janela intermediária captura o entorno imediato da lacuna de forma mais estável.
Uma janela maior aproxima um contexto regional do recorte.
Isso significa que a mesma região pode exibir padrões de entropia distintos conforme a escala considerada.
Logo, não há motivo para assumir, a priori, que a relação entre entropia e qualidade seja a mesma em todas as escalas.

Do ponto de vista metodológico, os algoritmos avaliados no artigo pertencem a dois regimes de modelagem distintos.
Os métodos clássicos dependem de hipóteses explícitas sobre suavidade, proximidade espacial, correlação, esparsidade em bases transformadas ou autossimilaridade da própria imagem.
Esses métodos costumam ser mais interpretáveis e, muitas vezes, independem de treinamento supervisionado.
Os modelos de aprendizado profundo, por sua vez, aprendem uma função de reconstrução a partir de pares de entrada degradada e alvo limpo.
Eles podem capturar regularidades complexas do domínio, mas também dependem de escolha arquitetural, desenho de treinamento e qualidade do conjunto de dados.

Uma distinção teórica importante para este trabalho é a diferença entre qualidade absoluta e robustez.
Qualidade absoluta diz respeito ao quão próxima a reconstrução está da referência em determinada condição experimental.
Robustez diz respeito ao quanto esse desempenho se mantém quando a condição se torna mais adversa, por exemplo com maior ruído ou maior complexidade espacial.
Essas duas dimensões podem caminhar juntas, mas não são equivalentes.
Um método pode liderar em qualidade e ainda assim sofrer degradação mais acentuada quando o cenário piora.

Também é importante distinguir comparação interna de comparação cruzada.
Comparação interna ocorre quando métodos submetidos ao mesmo contexto experimental são comparados entre si, permitindo conclusões fortes sobre sua hierarquia relativa.
Comparação cruzada ocorre quando se tenta confrontar métodos avaliados em regimes diferentes.
Neste artigo, as comparações internas são o foco principal.
As comparações cruzadas entre blocos clássico e profundo aparecem apenas como contexto interpretativo amplo, nunca como base para um ranking universal.

Essa escolha está diretamente relacionada à validade científica da análise.
Quando resultados de famílias metodológicas distintas são agregados sem ressalvas, o leitor pode confundir diferenças de domínio experimental com diferenças algorítmicas intrínsecas.
Ao separar os blocos, o estudo preserva a comparabilidade onde ela é mais forte e evita conclusões que excedam o alcance real dos dados.

Em síntese, a fundamentação deste trabalho pode ser condensada em uma ideia central.
O desempenho em preenchimento de lacunas emerge da interação entre método, estrutura local da cena, ruído no contexto observado e custo computacional.
A entropia local torna essa interação mensurável.
A separação entre métodos clássicos e aprendizado profundo torna sua interpretação mais rigorosa.
